% Suggested final budget: 1.5--2.0 pages, including plots and tables.
\section{Results}
\label{sec:results}
\subsection{Six-task baselines}
\begin{table}[htbp]
  \centering
  \caption{Placeholder for six-task DP baselines on held-out test initial
  conditions. Once and End denote \metriconce{} and \metricend{}.
  Final rates should be training-seed mean $\pm$ sample standard deviation.
  Identify the backbone, checkpoint, and any task-specific settings.}
  \label{tab:baselines}
  \small
  \begin{tabularx}{\linewidth}{@{}Yccccc@{}}
    \toprule
    Task & $N$ & Seeds & \shortstack{Episodes\\per seed} & Once (\%) & End (\%) \\
    \midrule
    PickCube & \pending & \pending & \pending & \pending & \pending \\
    StackCube & \pending & \pending & \pending & \pending & \pending \\
    PushCube & \pending & \pending & \pending & \pending & \pending \\
    PullCube & \pending & \pending & \pending & \pending & \pending \\
    PegInsertionSide & \pending & \pending & \pending & \pending & \pending \\
    PlugCharger & \pending & \pending & \pending & \pending & \pending \\
    \bottomrule
  \end{tabularx}
\end{table}

\todo{Summarize the verified per-task baseline results in
\cref{tab:baselines}. Report the evaluated configuration and episodes per seed.
Every value must be traceable to raw evaluation output. TBD denotes unavailable
data; it must never be interpreted as a measured zero.}

\subsection{Effect of demonstration count}
\begin{figure}[htbp]
  \centering
  \figplaceholder[45mm]{Six task panels: held-out success versus $\log_2 N$.
  Label the evaluated split, seed counts, and meaning of error bars.}
  \caption{Placeholder for the data-efficiency study. Final error bars should
  represent the stated seed-level variability or uncertainty measure.}
  \label{fig:data-efficiency}
\end{figure}
\todo{Report observed trends, adjacent-$N$ effect sizes and uncertainty, and
exceptions. Describe saturation only when supported. Do not infer improvement
from training loss alone or interpolate a meaningful data requirement when
success remains near zero.}

\subsection{Backbone success and efficiency}
\begin{table}[htbp]
  \centering
  \caption{Placeholder for a backbone comparison on one task and one dataset
  size. Specify task, $N$, training seed count, evaluation episodes, GPU model,
  and whether inference timing measures a denoising step or a full action chunk.
  Repeat task blocks as needed; do not average incompatible hardware timings.}
  \label{tab:backbones}
  \small
  \begin{tabularx}{\linewidth}{@{}Yccccc@{}}
    \toprule
    Backbone & \shortstack{Once\\(\%)} & \shortstack{Params\\(M)} &
    \shortstack{Train\\(ms/step)} & \shortstack{VRAM\\(GB)} &
    \shortstack{Infer\\(ms)} \\
    \midrule
    UNet & \pending & \pending & \pending & \pending & \pending \\
    Transformer & \pending & \pending & \pending & \pending & \pending \\
    MLP & \pending & \pending & \pending & \pending & \pending \\
    \bottomrule
  \end{tabularx}
\end{table}

\todo{Discuss supported differences and trade-offs in
\cref{tab:backbones}. State the task and $N$ for each block. Add a
success--inference-time plot only if it provides evidence beyond the table.
If this study is incomplete, narrow the scope or remove the subsection.}
\FloatBarrier
